%%%%%%%%%%%%%%%%%%%%%%%%%%%%%%%%%%%%%%%%%%%%%%%%%%%%%%%%%%%%%%%%%%%%%%
% Lớp 11 | Đường thẳng, mặt phẳng. Quan hệ song song trong | Phép chiếu song song
% Dạng 1H4?6-4: Xác định phương chiếu
% Số câu: 2  (sinh tự động bởi tach_cau.py; từ nay thầy sửa trực tiếp file này)
%%%%%%%%%%%%%%%%%%%%%%%%%%%%%%%%%%%%%%%%%%%%%%%%%%%%%%%%%%%%%%%%%%%%%%

%%%%%%%%%%%%%%%% TRẮC NGHIỆM 4 PHƯƠNG ÁN %%%%%%%%%%%%%%%%

% nguồn: Ok_TracNghiemLop11_Lan1
\begin{ex}%[1H4N6-4]%[ma:874ab5f4]
	Xét phép chiếu song song  lên mặt phẳng $\left(P\right)$ theo phương $l$. Trong các sau mệnh đề nào đúng?
	\choice 
		{Hình chiếu song song của hai đường thẳng cắt nhau có thể song song  với nhau}
		{\True Hình chiếu song song của hai  đường thẳng chéo nhau có thể song song với nhau}
		{Hình chiếu song song của hai đường thẳng chéo nhau thì  song song với nhau}
		{Hình  chiếu song song của hai đường thẳng chéo nhau thì không thể song song với nhau}
	\loigiai{ 
		Hình chiếu song song của hai đường thẳng chéo nhau có thể song song với nhau; cũng có thể cắt nhau. Do đó mệnh đề đúng là \lq\lq Hình chiếu song song của hai đường  thẳng chéo nhau có thể song song với nhau\rq\rq.
	} 
\end{ex}


%%%%%%%%%%%%%%%% ĐÚNG/SAI %%%%%%%%%%%%%%%%

% nguồn: Ok_dungsai_lop11_Lan1
\begin{ex}%[1H4H6-4]%[ma:5dd1ffa5]
  Cho hình lăng trụ $ABC.A'B'C'$, $I$ và $I'$ lần lượt là trung điểm của đoạn $AB$ và $A'B'$.
 \choiceTF 
  {\True $AI'  \parallel IB$}
  {Hình chiếu song song của $I$ trên mặt phẳng $(A'B'C')$ phương $A'I$ là điểm  $C'$}
  {\True Trong mặt phẳng $(A'B'C')$,  vẽ hình bình hành $A'C'MI'$. Suy ra $ACMI'$ là hình bình hành}
  {\True $M$ là hình chiếu song song  của $C$ theo phương $AI'$ trên mặt phẳng $(A'B'C')$}
 \loigiai{\begin{center}
  \begin{tikzpicture}[line join=round,line cap=round,>=stealth,thick]
    \coordinate (A)  at (-2,3);
     \coordinate (B) at (3,3);
    \coordinate (C)  at (0,1.6);
    \coordinate  (A') at (-2.6,-2);
    \coordinate (B') at (2.4,-2); 
    \coordinate (C')  at (-0.5,-3.4);
    \coordinate  (I) at ($(A)!0.5!(B)$);
    \coordinate  (I') at ($(A')!0.5!(B')$);
    \coordinate (M) at  (2,-3.4);
    \draw[dashed](A')--(B') 
      (A)--(I')
     (I)--(B') 
     (I')--(M) 
     (B')--(C'); 
    \draw (A)--(B) 
      (A)--(C)
     (C)--(B) 
      (A')--(C')
      (C')--(M)
      (C)--(C')
     (A)--(A') 
     (B)--(B') 
     (C)--(M); 
     \foreach \i in {A,B,C,A',B',C',I,I',M}{\draw[fill=](\i) circle (1pt);}
    \foreach \i in  {A,I,B}{\draw (\i) node[scale=0.9,above]{$\i$};}
    \foreach  \i in {A',C',C}{\draw (\i) node[scale=0.9,left]{$\i$};}
    \foreach \i  in {B',M}{\draw (\i) node[scale=0.9,right]{$\i$};}
    \foreach  \i in {I'}{\draw (\i) node[scale=0.9,above right]{$\i$};}
   \end{tikzpicture}
 \end{center}
 \begin{itemchoice}
  \itemch \textbf{Đúng}. Ta có $\heva{&AI \parallel B'I'\\ &AI=B'I'=\dfrac{AB}{2}} \Rightarrow AIB'I'$ là hình bình hành, do đó $AI' \parallel IB'$. 
  \itemch \textbf{sai}. Vậy hình chiếu song song của $I$ trên mặt phẳng $(A'B'C')$ phương $A'I$  là điểm $B'$.
  \itemch \textbf{Đúng}. Trong  mặt phẳng $(A'B'C')$, vẽ hình bình hành $A' C'M I'$.\\
  Vì $\heva{&MI' \parallel AC',\,MI'=AC' \\& AC' \parallel AC,\,  AC'=AC}\Rightarrow MI' \parallel AC,\, MI'=AC$.\\
  Suy ra  $ACMI'$ là hình bình hành.
  \itemch \textbf{Đúng}.  Vì vậy $AI'\parallel CM$, mà $M \in (A'B'C')$ nên $M$ chính là hình chiếu song song của $C$ theo phương $AI'$ trên mặt phẳng $(A'B'C')$.
 \end{itemchoice}
} 
\end{ex}
